%------------------------
% Resume Template
% Author : Harsh
% Tailored for: Airbus Operations GmbH -- "F&DT Lead Engineer for Fuselage
% Primary Structure (d/f/m)", Hamburg/Bremen (posting JR10428090)
%
% OPEN ITEMS before sending (see project memory cv_verified_metrics.md):
%  - [x] PhD expected submission/defence date in the header -- fill in real
%        month/year (this was never captured in any earlier CV in this repo;
%        ask Harsh directly, he will know his own timeline).
%  - Honest gap, not filled with a placeholder because there is no true
%        number to put here: Harsh has no formal line-management / "employees
%        report to me" leadership experience. His leadership evidence is
%        technical: instructing/supervising 60+ Master's students and
%        sign-off-style review of externally and student-produced analysis.
%        The cover letter carries an explicit transparency paragraph on this.
%  - Honest gap: no repair-scheme design/validation experience (posting asks
%        for "develop and validate effective repair schemes"). Not claimed
%        anywhere in this CV. Flagged in the cover letter too.
%  - Hyperworks is a "nice to have" in the posting -- Harsh has not used it
%        (Nastran/Patran + ABAQUS only) and it is deliberately not listed.
%------------------------

\documentclass[a4paper,20pt]{article}

\usepackage{latexsym}
\usepackage[empty]{fullpage}
\usepackage{titlesec}
\usepackage{marvosym}
\usepackage[usenames,dvipsnames]{color}
\usepackage{verbatim}
\usepackage{enumitem}
\usepackage[pdftex]{hyperref}
\usepackage{fancyhdr}
\usepackage{graphicx}
\usepackage{tikz}
\usepackage{blindtext}
\usepackage[document]{ragged2e}
\usepackage{xcolor}
\usepackage{hhline}
\usepackage{hyperref}
\usepackage{xcolor}

\hypersetup{
    colorlinks=true,
    linkcolor=blue,
    urlcolor=blue,
    citecolor=blue
}
\usepackage{tikz}

% Helper command to draw skill circles
\newcommand{\skilllevel}[1]{
  \begin{tikzpicture}[baseline=-0.5ex]
    \foreach \i in {1,...,5} {
      \ifnum\i>#1
        \draw[fill=white] (\i*0.25,0) circle (0.08);
      \else
        \draw[fill=black] (\i*0.25,0) circle (0.08);
      \fi
    }
  \end{tikzpicture}
}


\pagestyle{fancy}
\fancyhf{} % clear all header and footer fields
\fancyfoot{}
\renewcommand{\headrulewidth}{0pt}
\renewcommand{\footrulewidth}{0pt}

% Adjust margins
\addtolength{\oddsidemargin}{-0.530in}
\addtolength{\evensidemargin}{-0.375in}
\addtolength{\textwidth}{1in}
\addtolength{\topmargin}{-.45in}
\addtolength{\textheight}{1in}

\urlstyle{rm}

\raggedbottom
\raggedright
\setlength{\tabcolsep}{0in}

% Sections formatting
\titleformat{\section}{
  \vspace{-10pt}\scshape\raggedright\large
}{}{0em}{}[\color{blue}\titlerule \vspace{-6pt}]

%-------------------------
% Custom commands
\newcommand{\resumeItem}[2]{
  \item\small{
    \textbf{#1}{: #2 \vspace{-2pt}}
  }
}

\newcommand{\resumeItemWithoutTitle}[1]{
  \item\small{
    {\vspace{-2pt}}
  }
}

\newcommand{\resumeSubheading}[4]{
  \vspace{-1pt}\item
    \begin{tabular*}{0.97\textwidth}{l@{\extracolsep{\fill}}r}
      \textbf{#1} & #2 \\
      \textit{#3} & \textit{#4} \\
    \end{tabular*}\vspace{-5pt}
}

% Add this to your preamble
\usepackage{array}
\newenvironment{awlist}{% 
    \begin{tabular}{>{\bfseries}l @{\hspace{1em}} p{13cm}}
}{%
    \end{tabular}
}


\newcommand{\resumeSubItem}[2]{\resumeItem{#1}{#2}\vspace{-3pt}}

\renewcommand{\labelitemii}{$\circ$}

\newcommand{\resumeSubHeadingListStart}{\begin{itemize}[leftmargin=*]}
\newcommand{\resumeSubHeadingListEnd}{\end{itemize}}
\newcommand{\resumeItemListStart}{\begin{itemize}}
\newcommand{\resumeItemListEnd}{\end{itemize}\vspace{-5pt}}
%\newcommand{\resumeSubHeadingListStart}{\begin{itemize}[label={}, leftmargin=0pt]}
%\newcommand{\resumeSubHeadingListEnd}{\end{itemize}}

\usepackage{fontawesome5}
\usepackage{hyperref}
\hypersetup{colorlinks=true, urlcolor=blue}

%\usepackage{tgheros}
%\renewcommand{\familydefault}{\sfdefault}

\usepackage{charter}
%\usepackage{fourier}
%\usepackage{tgpagella}
%\usepackage{lmodern}
%\usepackage{tgheros}
%\usepackage{newtxtext}

%-----------------------------
%%%%%%  CV STARTS HERE  %%%%%%

\begin{document}




%----------HEADING-----------------

\begin{center}
    \textbf{{\LARGE Harsh Bordekar}} \\
    \textbf{F\&DT Lead Engineer for Fuselage Primary Structure (d/f/m)} \\
    \textbf{PhD Candidate at DLR \& TU Braunschweig, expected \textbf{[x]}} \\
    \textbf{Fatigue \& Damage Tolerance: Metallic \& Composite Structures, Model V\&V, Technical Review}
    
    \vspace{0.5em}
    
    \small
    \begin{tabular}{*{3}{l@{\hspace{2.5em}}}}
        \faIcon[regular]{envelope} \href{mailto:harshbordekar@gmail.com}{harshbordekar@gmail.com} & 
        \faIcon{linkedin} \href{https://www.linkedin.com/in/harsh-s-bordekar/}{Harsh's LinkedIn} & 
%        \faIcon{github} \href{https://github.com/gagankaushikmanyam}{Gagan's Github} \\
        \faIcon{mobile-alt} +49 (0) 17673888310 &
        \faIcon{map-marker-alt} Braunschweig, Germany & 
        %\faIcon{home} Niedersachsen (Residence) & 
        \faIcon{id-card} Permanent Resident - Germany
    \end{tabular}
\end{center}



\vspace{5pt}

\section{\color{blue}Profile}
Structural analysis and simulation engineer with 7+ years of experience in aerospace structural integrity, spanning fatigue and damage tolerance (F\&DT) of metallic structures and progressive damage of composites. Experience includes deriving damage-tolerance philosophies and allowable defect/damage limits for a target fatigue life using small-defect fatigue models (Murakami, El-Haddad) and certification-aligned failure criteria, together with high-fidelity linear and non-linear FEM, model verification \& validation against measured data, and stochastic/uncertainty analysis for certification-oriented sizing. Currently a PhD Candidate at DLR \& TU Braunschweig, developing a certification-aligned damage/leakage failure criterion that translates multiscale damage models into structural sizing inputs consistent with regulatory boundary conditions, and providing the stress justification required for structural approval. Regularly provides sign-off-style technical review of externally and student-produced structural analysis work, defines the technical scope for that work, and communicates findings to design and project teams. Familiar with analytical stress methods (Niu, Bruhn) alongside FE solvers (Nastran/Patran, ABAQUS), with a strong record of process automation using Python and HPC pipelines and an active interest in applying AI to improve analysis workflows. Works across interfaces with partner institutes, students, and stakeholders from varied technical and cultural backgrounds.

\vspace{0.2cm}
\textbf{Core Competencies:} Fatigue \& Damage Tolerance (F\&DT), Damage Tolerance Philosophy \& Allowable Limits, Fatigue Life \& Allowable Defect Size (Murakami, El-Haddad), Metallic \& Composite Structural Analysis, Linear \& Non-linear FEM (Nastran/Patran, ABAQUS), Analytical Stress Methods (Niu, Bruhn), Verification \& Validation, Certification-Oriented Sizing, Stress Approval of Analyses \& Technical Review, Simulation \& Reporting Automation, Uncertainty Quantification (UQ), Python Programming, HPC (SLURM)



\section{\color{blue}Professional Experience}


\resumeSubheading
    {Structural Analysis / Simulation Engineer}{DLR \& TU Braunschweig}
    {Damage tolerance philosophy, technical review \& simulation automation}{Feb $2021$ - Present}
    \resumeItemListStart
    \item (PhD) Derived a certification-aligned leakage failure criterion (leakage envelope) that maps combined strain/load states to leakage onset, defining an allowable-damage-limit philosophy that bridges progressive damage models and system-level tank simulations and translates multiscale FEM outputs into structural sizing inputs consistent with regulatory boundary conditions, providing the failure-criterion definition and stress justification required for structural approval.
\\
\textbf{Tools: ABAQUS, Python, FEM post-processing using C++}

\item Verified and technically reviewed externally and student-produced analysis models, meshes, and simulation results against defined methods and acceptance criteria, providing sign-off-style checks on structural analysis data and documentation.
\\
\textbf{Tools: ABAQUS, Nastran, Python}

\item Instructed 60+ Master's students in FEM theory and ABAQUS-based CFRP modelling, covering Continuum Damage Mechanics (CDM), Progressive Damage Analysis (PDA), and Cohesive Zone Modelling (CZM), demonstrating both deep technical expertise and the ability to lead and communicate complex simulation workflows.
\\
\textbf{Tools: ABAQUS, Automation using Python}

\item (PhD) Developed a probabilistic, physics-based multiscale framework for leakage onset prediction in CFRP hydrogen storage tanks, coupling microstructure-informed mesoscale damage models with stochastic spatial defect distributions to propagate failure across ply, laminate, and structural scales and estimate probability of failure.
\\
\textbf{Tools: Git, Docker, Python, ABAQUS, FEM, HPC/SLURM}

\item Performed uncertainty quantification to generate leakage probability vs. strain relationships, enabling risk-informed design decisions for next-generation CFRP pressure vessels under hydrogen storage conditions.
\\
\textbf{Tools: Python (scipy, numpy), Monte Carlo methods, ABAQUS}

\item Developed an Agentic AI application for microstructure-informed property prediction, accepting natural language queries and microstructure images as inputs to predict homogenized/effective material properties, powered by a physics-based CNN as a custom tool and a hierarchical RAG pipeline for context-aware retrieval of materials knowledge, enabling on-demand computational homogenization without full-scale FEM runs.
\\
\textbf{Tools: Python, Physics-based CNN, Hierarchical RAG, LLM agent framework, image processing}
\item Evaluated thermal protection effectiveness for the Callisto reusable rocket program's metal upper housing bay, executing large-scale thermal and structural simulations on HPC infrastructure.
\\
\textbf{Tools: ABAQUS, SLURM (HPC), Python, Nastran, Patran}

\resumeItemListEnd
	

  \vspace{10pt}


\resumeSubheading
    {Research Intern, Metallic Fatigue \& Damage Tolerance}{DLR-Braunschweig}
    {Fatigue-life, allowable defect size \& damage tolerance qualification}{Nov $2019$ - Jan $2021$}
\resumeItemListStart
    \item Assessed fatigue and defect tolerance of aerospace-grade titanium (metallic) components, deriving allowable defect sizes for a target fatigue life using small-defect fatigue models (Murakami, El-Haddad), linking as-built defect populations to structural fatigue substantiation.
    \\
    \textbf{Tools: Small-defect fatigue models (Murakami, El-Haddad), Python (NumPy/SciPy)}

    \item Coupled defect measurements with numerical simulation to establish a defect-tolerance qualification methodology, translating fracture-mechanics and fatigue criteria into allowable-defect acceptance limits for additively manufactured metallic parts.
    \\
    \textbf{Tools: FEM (ABAQUS), fracture mechanics, Python}

    \item Developed a CT-based, explainable-AI defect-detection framework providing interpretable confidence outputs for engineering sign-off, automating classification of internal defects feeding the fatigue assessment.
    \\
    \textbf{Tools: Python (scikit-learn / TensorFlow, SHAP / LIME / Grad-CAM), CT image processing}


\resumeItemListEnd

  \vspace{10pt}



	
\resumeSubheading
    {Student Research Assistant}{Leibniz University Hannover}
    {Institute of Static and Dynamic}{Sep $2018$ - Oct $2019$}
    
    \resumeItemListStart
    

    \item Developed an ML framework for automated detection and quantification of fiber undulations in unidirectional CFRP laminates, generating spatially resolved defect metrics as direct inputs for probabilistic structural simulation models to assess manufacturing-induced variability on mechanical performance.
    \\
    \textbf{Tools: Matlab, Python (scikit-learn / TensorFlow), image processing, NumPy/SciPy}
    
\resumeItemListEnd
\vspace{10pt}

    \resumeSubheading
    {Design Engineer }{Gujarat Industrial Development Corporation}
    {AK. Alloys}{Mar $2017$ - Mar $2018$}
    \resumeItemListStart
   \item Designed sand cast components with manufacturing tolerances and process constraints, translating functional requirements into production-ready geometries compliant with sand casting process limitations, draft angles, parting line constraints, and dimensional quality standards.
       \\
    \textbf{Tools: CREO, SolidWorks, AutoCAD}

\item Managed witness casting and end-to-end quality assurance processes, ensuring dimensional conformance and material integrity of cast components against engineering specifications and customer acceptance criteria.
\\
\textbf{Tools: CREO, SolidWorks, AutoCAD}

\item Planned and optimized sand casting production schedules based on capacity, material availability, and delivery requirements, improving throughput alignment with manufacturing targets.
\textbf{Tools: SolidWorks, AutoCAD, ERP / production planning tools}
\resumeItemListEnd
        
		



%-------------SKILLS SUMMARY------------------
\renewcommand{\arraystretch}{1.3}
\section{\color{blue}Skills Summary}
\begin{tabular}{|p{8cm}p{10.5cm}|}

\hline
\textbf{Fatigue \& Damage Tolerance} &
Fatigue-life assessment \(\vert\) Allowable defect size (Murakami, El-Haddad) \(\vert\) Defect tolerance \(\vert\) Damage tolerance \& certification concepts \\

\hline
\textbf{Stress Analysis \& FEM} &
FE solvers (Nastran/Patran, ABAQUS) \(\vert\) Analytical methods (Niu, Bruhn) \(\vert\) Linear \& non-linear FEM \(\vert\) Verification \& Validation \\

\hline
\textbf{Numerical Simulations} &
Anisotropic material modeling and FEM  \(\vert\) Fracture Mechanics  \(\vert\) Progressive damage analysis \(\vert\) Multiscale Modeling \\

\hline
\textbf{ Programming Languages} &
Python  \(\vert\) C++ \(\vert\) FORTRAN \(\vert\) MATLAB \\
\hline

\textbf{ Data Analysis \& Scientific Computing} &
NumPy \(\vert\) Pandas \(\vert\) Matplotlib \(\vert\) Seaborn \(\vert\) Plotly \(\vert\) SALib \(\vert\) SciPy \\
\hline

\textbf{ ML \& AI} &
Scikit-learn \(\vert\) TensorFlow \(\vert\) PyTorch \\
\hline

\textbf{LLMs \& Generative AI} &
Hugging Face Transformers (Fine-tuning) \(\vert\) Prompt Engineering \(\vert\) \\ & Embeddings \(\vert\) RAG \(\vert\) LangChain \(\vert\) LangGraph \(\vert\) MCP \\

\hline

\textbf{ APIs \& Backend Development} &
RESTful APIs \(\vert\) FastAPI \(\vert\) JSON\\
\hline


\textbf{ Version Control \& Reproducibility} &
Git \(\vert\) Docker \\
\hline

\textbf{ Application Development} &
VS Code \(\vert\) JupyterLab \(\vert\) Google Colab \(\vert\) Streamlit \(\vert\) Gradio\\
\hline
\textbf{ Languages} &
English (C1) \(\vert\) German (B1) \\
\hline
\end{tabular}

\vspace{0.2cm}
%\small{\textit{Keywords: Machine Learning, Applied ML, Model Training, Model %Evaluation, Feature Engineering, Inference Pipelines, Reproducible ML}}
\vspace{5pt}

\newpage
%-----------Awards-----------------

%-----------EDUCATION-----------------
\section{\color{blue}Awards and Education}

\resumeSubheading
  {\textcolor{green!50!black}{\faAward} \textbf{2nd Prize: Best Paper Award}}
  {DLR (German Aerospace Center)}{}{2026}
\resumeSubheading
  {\textcolor{green!50!black}{\faAward} \textbf{3rd Prize: Best Paper Award}}
  {DLR (German Aerospace Center)}{}{2024}
\vspace{4pt}

\resumeSubheading
  {\textcolor{green!50!black}{M.Sc.} Computational Methods in Engineering}
  {Leibniz University Hannover, Germany}{}{April 2018 -- Dec 2020}
  \vspace{2pt}
\textbf{Focus:} Numerical Methods, Scientific Machine Learning, Deep Learning
\vspace{4pt}
\section{\color{blue}Selected Publications}

\begin{enumerate}

  \item \textbf{Bordekar, H.}, et al.
  ``Microstructure-informed Probabilistic Multiscale Modelling of CFRP via Machine-Learning-Derived Stochastic Volume Elements.''
  \textit{Manuscript in preparation}, 2026.

  \item \textbf{Bordekar, H.}, et al.
  ``MicroStructAgent: A Physics-Informed Agentic AI Framework for Computational Homogenization and Quality Assurance of Fiber Composite Laminates.''
  \textit{Manuscript in preparation}, 2026.

  \item \textbf{Bordekar, H.}, et al.
  ``Microstructure Characterization of Fiber Composite Laminate Using eXplainable Artificial Intelligence.''
  \textit{Journal of Composite Materials}, vol. 60, no. 1, 2026, pp. 3--25.
  \href{https://journals.sagepub.com/doi/pdf/10.1177/00219983251345559}{\textcolor{blue}{[Link]}}

  \item \textbf{Bordekar, H.}, et al.
  ``eXplainable Artificial Intelligence for Automatic Defect Detection in Additively Manufactured Parts Using CT Scan Analysis.''
  \textit{Journal of Intelligent Manufacturing}, vol. 36, no. 2, 2025, pp. 957--974.
  \href{https://link.springer.com/article/10.1007/s10845-023-02272-4}{\textcolor{blue}{[Link]}}

\end{enumerate}
\vspace{2pt}

\section{\color{blue}References}

\begin{tabular}{@{} p{6cm} p{10cm}@{}}
\textbf{\href{https://www.dlr.de/de/sy/ueber-uns/abteilungen-dlr-sy/dlr-sy-funktionsleichtbau}{Prof. Dr. Christian Huehne}} & Abteilungsleiter Funktionsleichtbau. Deutsches Zentrum für Luft- und Raumfahrt. Institut für Systemleichtbau \\
\textbf{\href{https://www.tu-braunschweig.de/ima/team/head-of-institute/prof-dr-ing-oliver-voelkerink}{Prof. Dr. Oliver Voelkerink}} & Head of Institute., IMA, TU-Braunschweig\\
\textbf{\href{https://www.linkedin.com/in/nicola-cersullo-41772b131 }{
Dr. Nicola Cersullo}} & R\&T Project Manager @ Airbus Hamburg \\
\end{tabular}

\end{document}

